% GENERATED FILE. Edit canonical lessons, then run npm run generate:books.
% canonical-source-hash: fnv1a64:83b1fb7a34b6ee86
% canonical-lessons: FA-C356-riz-kardan, FA-C356-charkh-kardan, FA-C356-leh-kardan, FA-C356-taft-dadan, FA-C356-piyaz-dagh

\chapter{Chop, mince and fry}
\label{ch:fa-chop-mince-and-fry}
\hlchaptermodality{356}

\begin{chapteropening}
\textbf{By the end of this chapter:} \emph{I can say how food was chopped fine, minced, mashed and lightly fried, and talk about fried onion.}

Complete the last lesson of chapter 356: I can say how food was chopped fine, minced, mashed and lightly fried, and talk about fried onion.
\end{chapteropening}

\section[\texorpdfstring{\fa{ریز} \fa{کردن} (riz kardan)}{ریز کردن (riz kardan)}]{\fa{ریز} \fa{کردن} (riz kardan) — to chop finely, to dice}

\label{lesson:FA-C356-riz-kardan}



\begin{quote}
Before the new one: say the Persian for to sift, then the Persian for to soak.
\end{quote}

\subsection*{You'll want to know: \fa{ریز} \fa{کردن}}
\textbf{\fa{ریز} \fa{کردن}} (\emph{riz kardan}) — \textquotedblleft{}to chop finely, to dice\textquotedblright{}.

\begin{quote}
\fa{سبزی} \fa{ها} \fa{را} \fa{ریز} \fa{کرد}.
\end{quote}

\begin{quote}
\emph{Sabzi-hâ râ riz kard.}
\end{quote}

\begin{quote}
She chopped the herbs finely.
\end{quote}

\emph{Riz} \textquotedblleft{}small, fine\textquotedblright{}. Present \emph{riz mikonad}.

\begin{cousinweb}
\emph{riz} \textquotedblleft{}fine, small\textquotedblright{} is a native Persian word that EtymDB lists without a parent; \emph{kardan} \textquotedblleft{}to do\textquotedblright{} is inherited Persian.
\end{cousinweb}

\begin{grammarlens}[title={What you've built}]
One more word for telling, explaining and describing at B1.
\end{grammarlens}

\subsection*{Your turn}
\begin{itemize}
\raggedright
  \item \emph{Say it:} \emph{riz kardan}
  \item \emph{Say it:} \emph{riz kardan}, once more
  \item \emph{Say it:} a sentence of your own with \emph{riz kardan}
  \item \emph{Recall:} say the Persian for to sift, then the Persian for to soak, then say \emph{riz kardan} again
\end{itemize}

\begin{tcolorbox}[breakable,colback=teal!4,colframe=teal!35!black,title={Before you move on}]
What does \fa{ریز} \fa{کردن} mean? (\textquotedblleft{}To chop finely, to dice\textquotedblright{}.) Say it once more.
\end{tcolorbox}
\section[\texorpdfstring{\fa{چرخ} \fa{کردن} (charkh kardan)}{چرخ کردن (charkh kardan)}]{\fa{چرخ} \fa{کردن} (charkh kardan) — to mince, to grind (meat)}

\label{lesson:FA-C356-charkh-kardan}



\begin{quote}
Before the new one: say the Persian for to soak, then the Persian for to chop finely.
\end{quote}

\subsection*{You'll want to know: \fa{چرخ} \fa{کردن}}
\textbf{\fa{چرخ} \fa{کردن}} (\emph{charkh kardan}) — \textquotedblleft{}to mince, to grind (meat)\textquotedblright{}.

\begin{quote}
\fa{قصاب} \fa{گوشت} \fa{را} \fa{چرخ} \fa{کرد}.
\end{quote}

\begin{quote}
\emph{Qassâb gusht râ charkh kard.}
\end{quote}

\begin{quote}
The butcher minced the meat.
\end{quote}

\emph{Gusht-e charkh karde} is minced meat. \emph{Charkh} \textquotedblleft{}wheel\textquotedblright{}, for the mincer.

\begin{cousinweb}
\emph{charkh} \textquotedblleft{}wheel\textquotedblright{} is inherited: EtymDB traces it through Middle Persian to the Indo-European word for a wheel, \emph{k\textsuperscript{w}ék\textsuperscript{w}los}. \emph{kardan} \textquotedblleft{}to do\textquotedblright{} is inherited too.
\end{cousinweb}

\begin{grammarlens}[title={What you've built}]
One more word for telling, explaining and describing at B1.
\end{grammarlens}

\subsection*{Your turn}
\begin{itemize}
\raggedright
  \item \emph{Say it:} \emph{charkh kardan}
  \item \emph{Say it:} \emph{charkh kardan}, once more
  \item \emph{Say it:} a sentence of your own with \emph{charkh kardan}
  \item \emph{Recall:} say the Persian for to soak, then the Persian for to chop finely, then say \emph{charkh kardan} again
\end{itemize}

\begin{tcolorbox}[breakable,colback=teal!4,colframe=teal!35!black,title={Before you move on}]
What does \fa{چرخ} \fa{کردن} mean? (\textquotedblleft{}To mince, to grind (meat)\textquotedblright{}.) Say it once more.
\end{tcolorbox}
\section[\texorpdfstring{\fa{له} \fa{کردن} (leh kardan)}{له کردن (leh kardan)}]{\fa{له} \fa{کردن} (leh kardan) — to mash, to crush}

\label{lesson:FA-C356-leh-kardan}



\begin{quote}
Before the new one: say the Persian for to chop finely, then the Persian for to mince.
\end{quote}

\subsection*{You'll want to know: \fa{له} \fa{کردن}}
\textbf{\fa{له} \fa{کردن}} (\emph{leh kardan}) — \textquotedblleft{}to mash, to crush\textquotedblright{}.

\begin{quote}
\fa{سیب} \fa{زمینی} \fa{ها} \fa{را} \fa{با} \fa{چنگال} \fa{له} \fa{کرد}.
\end{quote}

\begin{quote}
\emph{Sibzamini-hâ râ bâ changâl leh kard.}
\end{quote}

\begin{quote}
She mashed the potatoes with a fork.
\end{quote}

\emph{Leh} \textquotedblleft{}crushed, squashed\textquotedblright{}. \emph{Leh shodan} \textquotedblleft{}to get squashed\textquotedblright{}.

\begin{cousinweb}
\emph{leh} \textquotedblleft{}crushed\textquotedblright{} is a Persian word that EtymDB lists without a parent; \emph{kardan} \textquotedblleft{}to do\textquotedblright{} is inherited Persian.
\end{cousinweb}

\begin{grammarlens}[title={What you've built}]
One more word for telling, explaining and describing at B1.
\end{grammarlens}

\subsection*{Your turn}
\begin{itemize}
\raggedright
  \item \emph{Say it:} \emph{leh kardan}
  \item \emph{Say it:} \emph{leh kardan}, once more
  \item \emph{Say it:} a sentence of your own with \emph{leh kardan}
  \item \emph{Recall:} say the Persian for to chop finely, then the Persian for to mince, then say \emph{leh kardan} again
\end{itemize}

\begin{tcolorbox}[breakable,colback=teal!4,colframe=teal!35!black,title={Before you move on}]
What does \fa{له} \fa{کردن} mean? (\textquotedblleft{}To mash, to crush\textquotedblright{}.) Say it once more.
\end{tcolorbox}
\section[\texorpdfstring{\fa{تفت} \fa{دادن} (taft dâdan)}{تفت دادن (taft dâdan)}]{\fa{تفت} \fa{دادن} (taft dâdan) — to sauté, to fry lightly}

\label{lesson:FA-C356-taft-dadan}



\begin{quote}
Before the new one: say the Persian for to mince, then the Persian for to mash.
\end{quote}

\subsection*{You'll want to know: \fa{تفت} \fa{دادن}}
\textbf{\fa{تفت} \fa{دادن}} (\emph{taft dâdan}) — \textquotedblleft{}to sauté, to fry lightly\textquotedblright{}.

\begin{quote}
\fa{پیاز} \fa{را} \fa{در} \fa{کره} \fa{تفت} \fa{داد}.
\end{quote}

\begin{quote}
\emph{Piyâz râ dar kare taft dâd.}
\end{quote}

\begin{quote}
She lightly fried the onion in butter.
\end{quote}

Quick frying while stirring. Present \emph{taft midahad}.

\begin{cousinweb}
\emph{taft} \textquotedblleft{}heat\textquotedblright{} is the past stem of the native verb \emph{taftan} \textquotedblleft{}to heat, to glow\textquotedblright{}, and \emph{dâdan} \textquotedblleft{}to give\textquotedblright{} is inherited Persian: to give it heat.
\end{cousinweb}

\begin{grammarlens}[title={What you've built}]
One more word for telling, explaining and describing at B1.
\end{grammarlens}

\subsection*{Your turn}
\begin{itemize}
\raggedright
  \item \emph{Say it:} \emph{taft dâdan}
  \item \emph{Say it:} \emph{taft dâdan}, once more
  \item \emph{Say it:} a sentence of your own with \emph{taft dâdan}
  \item \emph{Recall:} say the Persian for to mince, then the Persian for to mash, then say \emph{taft dâdan} again
\end{itemize}

\begin{tcolorbox}[breakable,colback=teal!4,colframe=teal!35!black,title={Before you move on}]
What does \fa{تفت} \fa{دادن} mean? (\textquotedblleft{}To sauté, to fry lightly\textquotedblright{}.) Say it once more.
\end{tcolorbox}
\section[\texorpdfstring{\fa{پیاز} \fa{داغ} (piyâz dâgh)}{پیاز داغ (piyâz dâgh)}]{\fa{پیاز} \fa{داغ} (piyâz dâgh) — fried onion}

\label{lesson:FA-C356-piyaz-dagh}



\begin{quote}
Before the new one: say the Persian for to mash, then the Persian for to sauté.
\end{quote}

\subsection*{You'll want to know: \fa{پیاز} \fa{داغ}}
\textbf{\fa{پیاز} \fa{داغ}} (\emph{piyâz dâgh}) — \textquotedblleft{}fried onion\textquotedblright{}.

\begin{quote}
\fa{روی} \fa{آش} \fa{پیاز} \fa{داغ} \fa{ریخت}.
\end{quote}

\begin{quote}
\emph{Ru-ye âsh piyâz dâgh rikht.}
\end{quote}

\begin{quote}
He spooned fried onion over the soup.
\end{quote}

Thinly sliced onion fried until golden, a topping for many dishes. \emph{Dâgh} \textquotedblleft{}hot\textquotedblright{}.

\begin{cousinweb}
\emph{piyâz} \textquotedblleft{}onion\textquotedblright{} is inherited from Middle Persian (EtymDB), and \emph{dâgh} \textquotedblleft{}hot\textquotedblright{} is native too; EtymDB takes it back to the Indo-European root \emph{d\textsuperscript{h}eg\textsuperscript{w}\textsuperscript{h}-} \textquotedblleft{}to burn\textquotedblright{}.
\end{cousinweb}

\begin{grammarlens}[title={What you've built}]
One more word for telling, explaining and describing at B1.
\end{grammarlens}

\subsection*{Your turn}
\begin{itemize}
\raggedright
  \item \emph{Say it:} \emph{piyâz dâgh}
  \item \emph{Say it:} \emph{piyâz dâgh}, once more
  \item \emph{Say it:} a sentence of your own with \emph{piyâz dâgh}
  \item \emph{Recall:} say the Persian for to mash, then the Persian for to sauté, then say \emph{piyâz dâgh} again
\end{itemize}

\begin{tcolorbox}[breakable,colback=teal!4,colframe=teal!35!black,title={Before you move on}]
What does \fa{پیاز} \fa{داغ} mean? (\textquotedblleft{}Fried onion\textquotedblright{}.) Say it once more.
\end{tcolorbox}
